\documentclass[11pt]{article}
\usepackage[margin=1in]{geometry}
\usepackage{amsmath,amssymb}
\title{Disproof of Conjecture 00000007669}
\author{earthking11}
\date{6 October 2026}
\begin{document}
\maketitle

Starting from $(a,b)$, write $A$ for the move $a\mapsto aq$ and $B$ for
$b\mapsto bq$.  At depth three the binary move histories are
\[
 AAA,AAB,ABA,ABB,BAA,BAB,BBA,BBB.
\]
Thus, if the phrase ``binary tree'' means that histories are its nodes, the
level-three node count is $2^3=8$.

There is a second possible reading: identify histories that produce the same
parameter pair.  But $A$ and $B$ commute, and a word containing $i$ copies of
$A$ and $j$ copies of $B$ ends at $(aq^i,bq^j)$.  The depth-three endpoint
states are exactly
\[
 (3,0),(2,1),(1,2),(0,3),
\]
where the pair records $(i,j)$.  Hence the quotient count is $4$.

The conjecture instead asserts
\[
 a_3=\operatorname{Fib}(3+2)=\operatorname{Fib}(5)=5.
\]
It therefore fails under either natural meaning of ``node'': $8\ne5$ before
identification and $4\ne5$ after identification.  This counterexample uses
only the two moves supplied in the definition and does not depend on any
analytic property of a Bailey identity.

The Lean file exhaustively lists the eight histories, maps them to exponent
pairs, removes duplicates, and kernel-checks the counts $8$, $4$, and the
Fibonacci value $5$.
\end{document}
